\section{Introduction}

\begin{figure*}
    \centering
    \includegraphics[width=\linewidth]{figures/science-settings.png}
    %\vspace{-2em} % crunch
    \caption{\textnormal{To find barriers and opportunities for AI adoption, we visited workplaces and performed on-site interviews with 12 \rev{scientific practitioners} in a range of lab and field settings such as \textbf{(A)} behavioral neuroscience labs with surgery and bench areas (P1, P5), \textbf{(B)} field site for a primate behavioral scientist (P8), \textbf{(C)} cellular biochemistry lab (P12).}}
    \Description{Four images of various scientific settings are labeled with (A), (B), (C), and (D). Image (A) shows a behavioral neuroscience lab bench with a microscope, tools, and scattered notes for soldering a neurobiology implant device. In (B), multiple scientists in cleanroom suits stand around computers and machinery in a nuclear fusion laboratory. In (C), a macaque in a dry landscape interacts with an apparatus made out of acrylic panels. Image (D) provides a closer look at a lab bench, featuring many reagent bottles and sticky notes on the shelves and a notebook on the table top.}
    %\vspace{-1em} % crunch
    \label{fig:science-setting}
\end{figure*}

As AI continues to advance, global investment toward AI for science and the automation of knowledge discovery have led to widely publicized, high-profile breakthroughs. In 2024, teams were awarded the Nobel Prize in Chemistry for AlphaFold, an AI system that uses deep learning to predict 3D protein structures~\cite{jumper2021highly}, along with the Nobel Prize in Physics for the theoretical foundations that led to deep neural networks~\cite{nobelprizephysics2024}. Other fields have seen similar innovations, from drug discovery~\cite{wallach2015atomnet, insilico2024phase2} to climate prediction~\cite{lam2023learning}. Tech companies have also recently invested in using AI for science, such as Google's \emph{AI Co-Scientist} to aid in hypothesis deliberation, literature review, and data synthesis~\cite{gottweis2025towards}.

However, most of these initiatives have centered on computation, simulation, data analysis, and other forms of knowledge work that take place ``at the desk.'' While AI systems are promising for the computational and simulation-heavy areas of science, this is only one aspect of scientific practice. Decades of STS (Science and Technology Studies) and HCI research have shown that discovery frequently involves material interaction~\cite{latour2013laboratory} and embodied physical labor in labs and field sites~\cite{cetina1999epistemic}. A lot of critical scientific labor happens ``beyond the desk'' in high-stakes and physically-demanding environments. For instance, a nuclear physicist must don a full-body cleanroom suit to enter a high-security lab space where they fabricate delicate fuel capsules, a task where fine-grained physical coordination is vital and \emph{a moment's lapse of attention can derail months of hard work and put others in danger}. A biochemist similarly labors under intense timing constraints, running between incubators, centrifuges, and a series of sterile workspaces to keep cell cultures alive while following a choreographed experimental protocol. A field scientist hikes through the forests of the Republic of Congo to study chimpanzee behavior, improvising tools on-the-go in the face of unpredictable terrain and limited infrastructure. These are all real stories from the 12 \rev{\emph{scientific practitioners}}\footnote{\rev{In this paper, we use the term \textbf{\textit{scientific practitioners}} to specifically refer to people engaged in the day-to-day practices that produce scientific knowledge, including experimental setup, data collection, troubleshooting, interpretation, and the embodied labor of laboratory and field environments. Although workers of all ages and experience levels can engage in such practitioner labor, oftentimes much of this labor is taken on by junior lab staff (e.g., graduate students, postdocs, lab technicians) while senior scientific staff (e.g., professors, industry lab PIs) may focus more on ``at the desk'' knowledge work such as writing grant proposals, framing research papers, and giving talks to academic and broader public audiences.}} we interviewed for this paper (see Figures~\ref{fig:teaser} and~\ref{fig:science-setting}), and they illustrate how scientific practice across a range of domains can require physical presence, quick judgment, and expert hand-eye coordination. This sort of \textit{``beyond the desk''} labor currently exceeds what can be codified or simulated on the computer since it involves intimate contact with the physical world.

Despite ``beyond the desk'' science contributing critical findings to research and innovation, little is known about how contemporary AI tools fit into these workflows. For instance, HCI research in AI often focuses on its capabilities for screen-bound knowledge workflows such as programming, data science, creative writing, or UI/UX/visual design~\cite{guo2025from, laban2024beyond,vaithilingam2022expectation, zheng2025towards, liu2023code,shi2023understanding, suh2024luminate}. However, scientific work is often constrained by physical realities and unpredictable environments, involves improvised non-repeatable procedures, and is shaped by the multi-sensory \emph{tacit knowledge}\footnote{\emph{Tacit knowledge} is hands-on, domain-specific knowledge that is hard to precisely articulate in words, and thus rarely written down~\cite{polanyi2009tacit}. In scientific settings this may involve a nuclear physicist knowing how to finely manipulate a fuel capsule `by feel' using their sensory intuition. By definition this \emph{unwritten} knowledge cannot be in the training sets for text-based AI systems like LLMs. One can imagine training AI with video data, but even those cannot capture senses like how some fuel capsules `feel' right or wrong when manipulated with precision handheld tools, and how to adjust on-the-fly.} of domain experts. Moreover, the stakes of error are high, ranging from wasted months of research effort and materials to putting humans in danger. These conditions may pose unique barriers to AI adoption, which led us to the following question that, to our knowledge, we are the first to raise: \textbf{How do \rev{scientific practitioners} who work in embodied, improvisational, and high-stakes lab and field environments perceive the relevance, limitations, and future potential of AI assistance?}

To address this question, we visited lab and field sites of 12 \rev{scientific practitioners} and conducted in-situ interviews in their workplaces. Our participants worked in laboratory and field environments across biology, neuroscience, animal cognition, nuclear physics, materials science, and field robotics. They had diverse backgrounds, from small university research labs to large industrial organizations, and a range of work experience levels \rev{ranging from junior university researchers to senior staff scientists.} Participants reflected on their current use of AI and articulated the stakes and limitations of AI assistance in their domains. At the end of each interview, they each engaged in a speculative design activity where they envisioned an imagined `ideal' future AI assistant.

Our study revealed three barriers that constrain AI adoption in scientific practice, shown in \fig{fig:teaser}: (A) experimental setups are too high-stakes to risk AI errors, (B) constrained environments make it hard to access AI, and (C) AI cannot match the tacit knowledge~\cite{polanyi2009tacit}, contextual judgment, and embodied physical skill of human experts. Thus, instead of seeking AI that ``does the science'' for them, they envisioned future tools that improve human memory, keep track of documentation, and help prevent costly errors due to lapses in human attention. They sketched a series of speculative designs for future AI systems, summarized in \fig{fig:newspecdesign}, that can 1) monitor task status, 2) organize lab-wide knowledge, 3) monitor scientists' health, 4) do field scouting, 5) do hands-on chores. Participants consistently favored passive, context-aware systems that could blend into existing workflows and support the conditions of human-scientific reasoning rather than AI that deprives them of the opportunity to \emph{``do the actual work [myself].''} (P11)

We hope our findings inform ongoing conversations within HCI and the broader Human-AI interaction community about how to design tools that respect the unique constraints of real-world scientific labor. In doing so, we also surface a broader opportunity to \emph{expand the design space of AI beyond raw productivity and automation, and toward supporting the intangible human conditions that make knowledge creation possible.} Our findings can inform AI design for other physically-based domains outside of science, such as medical practitioners, emergency first responders, artisanal craft workers, or field technicians working in demanding outdoor conditions. We make the following contributions to HCI:

\begin{itemize}
  \item The first study examining how \rev{scientific practitioners} perceive AI tools in embodied lab/field work `beyond the desk'
  \item Three sets of barriers that constrain AI use in materially grounded settings such as lab and field science (See \fig{fig:teaser})
  \item Five speculative design concepts reflecting \rev{scientific practitioners}’ future visions for AI support (See \fig{fig:newspecdesign})
\end{itemize}

